\definecolor{introblue}{RGB}{38, 87, 139}
\definecolor{introlight}{RGB}{238, 244, 251}
\definecolor{introline}{RGB}{70, 70, 70}
\definecolor{introgreen}{RGB}{234, 246, 236}

\resizebox{0.96\linewidth}{!}{%
\begin{tikzpicture}[
  font=\sffamily,
  >=Latex,
  box/.style={
    draw=introblue,
    rounded corners=5pt,
    line width=0.7pt,
    fill=introlight,
    align=center,
    text width=4.4cm,
    minimum height=2.2cm,
    inner sep=7pt,
    font=\sffamily\small,
  },
  resultbox/.style={
    draw=introblue,
    rounded corners=5pt,
    line width=0.7pt,
    fill=introgreen,
    align=center,
    text width=4.6cm,
    minimum height=2.2cm,
    inner sep=7pt,
    font=\sffamily\small,
  },
  arrow/.style={->, line width=0.7pt, draw=introline},
]

\node[box] (problem) at (0,0) {\textbf{Problem}\\[3pt]
Problematic LLM-generated\\
code review comments\\[4pt]
\footnotesize
Unsupported, irrelevant, incorrect,\\
non-actionable, or low-value feedback};

\node[box] (artifacts) at (6.2,0) {\textbf{Design artifacts}\\[3pt]
Operational taxonomy\\
Trade-off-aware evaluation schema\\
Controlled comparison protocol};

\node[resultbox] (evidence) at (12.5,0) {\textbf{Empirical evidence after execution}\\[3pt]
Failure-type results\\
Preservation and coverage results\\
Cost and escalation trade-offs};

\draw[arrow] (problem.east) -- (artifacts.west);
\draw[arrow] (artifacts.east) -- (evidence.west);

\node[font=\sffamily\footnotesize\itshape, text=introline, align=center]
  at (6.2,-2.15)
  {The paper moves from the review problem to reusable design artifacts, and then to empirical evidence once the study is executed.};

\end{tikzpicture}%
}
